% Generated draft: requires mathematical and source review.
\documentclass[11pt]{article}
\usepackage[margin=25mm]{geometry}
\usepackage{fontspec}
\setmainfont{Latin Modern Roman}
\usepackage{amsmath,amssymb,mathtools,unicode-math}
\setmathfont{Latin Modern Math}
\usepackage{enumitem,needspace,xurl,hyperref}
\setlength{\parindent}{0pt}
\setlength{\parskip}{5pt}
\setlength{\emergencystretch}{3em}
\newcommand{\N}{\mathbb{N}}
\newcommand{\Z}{\mathbb{Z}}
\newcommand{\Q}{\mathbb{Q}}
\newcommand{\R}{\mathbb{R}}
\newcommand{\C}{\mathbb{C}}
\newcommand{\F}{\mathbb{F}}
\newcommand{\A}{\mathbb{A}}
\newcommand{\PP}{\mathbb{P}}
\newcommand{\E}{\mathbb{E}}
\newcommand{\eps}{\varepsilon}
\newcommand{\dd}{\,\mathrm{d}}
\begin{document}
\section*{UnsolvedMath ID: 20002019}
\textbf{Source problem code:} catalogue entry
\subsection{Definitions and mathematical statement}
Problem 1c asks to ``Study variational problems arising from conformally invariant problems.'' This is a research agenda without a specified functional, domain, or conclusion. The aid places it under Thomas Branson's anti-conformal perturbations heading and identifies the following ``Michael Eastwood'' line as the next participant block, not attribution of this problem.

Its model is the Yamabe variational problem. Let \((M,g)\) be a closed Riemannian manifold of dimension \(n\ge3\), take the nonnegative Laplacian \(\Delta_g=-\operatorname{div}_g\nabla_g\), and set
\[
 a_n=\frac{4(n-1)}{n-2},\quad p=\frac{2n}{n-2},\quad
 L_g=a_n\Delta_g+R_g,
 \quad
 Y_g(u)=\frac{\int_M(a_n|\nabla u|^2+R_gu^2)\dd V_g}
 {(\int_M|u|^p\dd V_g)^{2/p}},\qquad u>0.
\]
For \(\widehat g=w^{4/(n-2)}g\), \(w>0\), the aid reports
\[
 L_{\widehat g}\phi=w^{-(p-1)}L_g(w\phi),\qquad
 Y_{\widehat g}(\phi)=Y_g(w\phi).
\]
The positive critical-point equation is \(L_gu=\kappa u^{p-1}\), with
\(\kappa=\int_MuL_gu\dd V_g/\int_Mu^p\dd V_g\).

Suppose \(\operatorname{Vol}_g(M)=1\) and \(R_g\equiv R\) is constant, so \(u=1\) is a normalized critical point. Write
\[
 E_g(u)=\int_M(a_n|\nabla u|^2+Ru^2)\dd V_g,
 \qquad C_g(u)=\int_Mu^p\dd V_g.
\]
For the constraint \(C_g=1\), the multiplier in \(E_g-\mu(C_g-1)\) is \(\mu=2R/p\). Its constrained Hessian on \(\int_Mv\dd V_g=0\) is reported as
\[
 \mathcal H(v,v)=2\int_M\bigl(a_n|\nabla v|^2-(p-2)Rv^2\bigr)\dd V_g.
\]
Since \((p-2)/a_n=1/(n-1)\), for connected \(M\) the exact infinitesimal nonnegative-stability criterion is
\[
 \lambda_1(\Delta_g)\ge\frac{R}{n-1},
\]
where \(\lambda_1\) is the first positive eigenvalue. Equality directions are the mean-zero eigenspace with eigenvalue \(R/(n-1)\). If \(R\le0\) and \(M\) is connected, the Hessian is strictly positive on nonzero mean-zero variations. These are Hessian assertions, not a claim of strict nonlinear local minimality in every degenerate case.

If \(M\) has \(b_0(M)\) components but only one total-volume constraint, componentwise constants satisfying the global mean-zero condition form a \(b_0(M)-1\)-dimensional subspace. At \(R=0\) these are exactly the additional zero modes; at \(R>0\) they are negative directions; at \(R<0\) they are positive. For nonconstant componentwise modes the same spectral threshold must be checked on each component. Thus connected stability cannot be transferred unchanged to a disconnected manifold.

The aid also invokes the Lichnerowicz--Obata sphere interpretation. In the explicit positive Einstein setting \(\operatorname{Ric}_g=(n-1)kg\), \(k>0\), one has \(R=n(n-1)k\) and the relevant bound is \(\lambda_1\ge nk\), with equality in the connected case characterized by the round sphere of curvature \(k\). This interpretation needs its Ricci-curvature hypothesis; constant scalar curvature alone does not supply that eigenvalue bound. The model illustrates one precisely normalized conformal variational problem rather than resolving the full agenda.
\subsection{Short English statement}
The source gives a broad agenda rather than one theorem. The aid uses the Yamabe quotient to make a model precise, including conformal covariance, the critical equation, and a coefficient-checked Hessian threshold. It retains the extra component-transfer modes when the manifold is disconnected and avoids attributing the problem to the following participant heading.
\subsection{Sources}
\begin{itemize}
\item Local UnsolvedMath registry, permanent ID 20002019, exact source code catalogue entry; canonical \texttt{aim-geometry-notes.json}, \texttt{notes[356]}, tag \texttt{problem}; AIM Workshop Problem Lists, Geometry, \emph{Conformal structure in geometry, analysis, and physics}, Problem 1c. Created and updated 14 August 2026.
\item Source URL \url{https://aimath.org/WWN/confstruct/confstruct.pdf}. The aid assigns Problem 1c to Thomas Branson's anti-conformal perturbations block and the following Michael Eastwood line to the next block. This is supplied contextual evidence; the PDF was not accessed.
\item The aid invokes the Yamabe quotient and Lichnerowicz--Obata sphere interpretation, without full bibliographic entries or theorem locations.
\item Attempt 1: \texttt{partial\_result}, problem status at run \texttt{open}, rigor label \texttt{proved}, novelty confidence \texttt{low}; classification \texttt{PARTIAL-PROGRESS}, registry status \texttt{partially\_solved}. Candidate: normalization diagnostic, multiplier, Hessian threshold, and disconnected corrections. These are dated source labels, not certification.
\item Ulam AI / UnsolvedMath Contributors; CC BY 4.0. Machine-generated research aid, not peer reviewed; independently verify before citation.
\end{itemize}
\end{document}
